% =====================================================================
%  Capítulo 2 — Fundamentação Teórica
% =====================================================================
\chapter{Fundamentação Teórica}
\label{cap:fundamentacao}

Este capítulo apresenta os conceitos necessários à compreensão do problema e
da solução desenvolvidos neste trabalho: vibração mecânica em máquinas
rotativas, aquisição e processamento digital de sinais, análise espectral por
FFT, catálogo de defeitos e assinaturas espectrais e o critério de
persistência seletiva \(R^3\). Os conceitos são conectados às diretrizes
técnicas que serviram de especificação para o sistema ARGUS.

\section{Vibração mecânica em máquinas rotativas}

A vibração mecânica é a resposta dinâmica de uma estrutura às forças
alternadas que atuam sobre ela. Em máquinas rotativas, essas forças têm origem
em fenômenos como desbalanceamento residual do rotor, desalinhamento de
eixos, folgas mecânicas, contato rotor-estator e defeitos de mancais
\cite{harris2002,randall2011}. Como cada fenômeno excita a estrutura em
frequências características, o sinal de vibração medido em um ponto da máquina
contém informação diagnóstica sobre o estado dos seus componentes.

A grandeza fundamental para a interpretação espectral é a \emph{frequência de
rotação} do eixo. Dada a rotação em rotações por minuto (RPM), a frequência
fundamental, em hertz, é
\begin{equation}
\label{eq:fr}
  f_r = \frac{RPM}{60}.
\end{equation}
As componentes espectrais de um defeito síncrono aparecem em múltiplos
inteiros ou fracionários dessa frequência fundamental. A componente de ordem
\(k\) (denotada \(kX\)) possui frequência
\begin{equation}
\label{eq:fk}
  f_{kX} = k\, f_r .
\end{equation}
Por exemplo, a componente de primeira ordem (\(1X\)) ocorre exatamente na
frequência de rotação; a de segunda ordem (\(2X\)) no dobro, e assim por
diante. A representação por \emph{ordens} tem a vantagem de ser independente
da rotação absoluta: um defeito caracterizado por \(2X\) dominante mantém essa
característica em qualquer regime de rotação, desde que as frequências sejam
recalculadas pela Equação~\eqref{eq:fr}. Essa é a convenção adotada tanto na
literatura de diagnóstico quanto na especificação do ARGUS, na qual todas as
frequências são derivadas da rotação e nunca armazenadas como valores
absolutos fixos.

A norma ISO~10816 estabelece critérios para a avaliação da severidade da
vibração medida em partes não rotativas de máquinas, e a qualificação de
analistas de monitoramento de condição por vibração é tratada pela série
ISO~18436 \cite{iso10816,iso18436}. A convenção de unidades adotada depende da
grandeza medida e da faixa de frequência de interesse, conforme a
Tabela~\ref{tab:unidades}, que reproduz a diretriz metrológica que originou o
projeto.

\begin{table}[htbp]
\centering
\caption{Parâmetro, unidade de engenharia, convenção e aplicação típica na análise de vibração.}
\label{tab:unidades}
\begin{tabular}{@{}llll@{}}
\toprule
\textbf{Parâmetro} & \textbf{Unidade} & \textbf{Convenção} & \textbf{Aplicação típica} \\
\midrule
Deslocamento & mils / $\mu$m & Pk-Pk  & Baixas frequências; mancais de deslizamento; órbitas \\
Velocidade   & in/s / mm/s  & Peak   & Severidade global; falhas mecânicas (1X, 2X, 3X) \\
Aceleração   & g's / m/s²   & RMS    & Altas frequências; rolamentos; engrenamento; cavitação \\
\bottomrule
\end{tabular}
\caption*{\footnotesize Fonte: diretrizes técnicas do projeto (docs/descricao.txt).}
\end{table}

\section{Aquisição de sinais de vibração}

A qualidade do diagnóstico depende diretamente da qualidade da aquisição. Dois
aspectos são centrais: o critério de amostragem e o condicionamento do sinal.

\subsection{Critério de Nyquist e resolução espectral}

Para digitalizar um sinal analógico sem perda de informação, a taxa de
amostragem \(f_s\) deve ser superior ao dobro da maior frequência de interesse
\(F_{max}\) (critério de Nyquist),
\begin{equation}
\label{eq:nyquist}
  f_s > 2\, F_{max},
\end{equation}
sendo o uso de filtros \emph{anti-aliasing} de hardware mandatório para impedir
que energias acima de \(f_s/2\) se \emph{aliasem} como componentes de baixa
frequência, mascarando inclusive sinais subsíncronos \cite{oppenheim2010}.
Quando \(F_{max}\) é definido pela ordem harmônica mais alta do defeito de
interesse, o critério da Equação~\eqref{eq:nyquist} determina o limite mínimo
de amostragem. No ARGUS, essa validação é feita de forma automática antes de
qualquer processamento (comportamento \emph{fail-fast}), rejeitando
configurações que violem o critério.

A \emph{resolução espectral} \(\Delta f\) --- a menor separação entre duas
frequências distinguíveis no espectro --- é determinada pela taxa de amostragem
e pelo número de amostras \(N\) do registro:
\begin{equation}
\label{eq:df}
  \Delta f = \frac{f_s}{N}.
\end{equation}
Uma resolução adequada é necessária para separar picos adjacentes e para
identificar bandas laterais (\emph{sidebands}), indicadores importantes de
severidade em falhas de rolamento e de modulação \cite{randall2011}.

\subsection{Janelamento}

Como a FFT pressupõe que o sinal seja periódico no registro analisado, o
truncamento de um sinal não periódico introduz vazamento espectral
(\emph{leakage}): a energia de uma componente se espalha para frequências
vizinhas. O janelamento pondera as amostras para mitigar esse efeito. A
Tabela~\ref{tab:janelas} resume os critérios de aplicação adotados como
referência metrológica.

\begin{table}[htbp]
\centering
\caption{Tipos de janela e critérios de aplicação no processamento espectral.}
\label{tab:janelas}
\begin{tabular}{@{}p{3cm}p{5.5cm}p{4.6cm}@{}}
\toprule
\textbf{Janela} & \textbf{Critério de aplicação} & \textbf{Impacto na precisão} \\
\midrule
Hanning   & Padrão para sinais contínuos em regime permanente & Equilíbrio entre seletividade e proteção contra dispersão \\
Flat Top  & Obrigatória para testes de aceitação e calibração & Prioriza a precisão absoluta da amplitude do pico \\
Rectangular & Restrita a sinais transientes que iniciam e terminam no registro & Eficaz apenas se o registro contém o transiente completo \\
\bottomrule
\end{tabular}
\caption*{\footnotesize Fonte: diretrizes técnicas do projeto (docs/descricao.txt).}
\end{table}

\section{Transformada de Fourier e análise espectral}

A análise espectral converte o sinal do domínio do tempo para o domínio da
frequência, revelando as componentes que compõem a vibração. Para um sinal
discreto \(x[n]\) com \(N\) amostras, a Transformada Discreta de Fourier (DFT)
é definida por
\begin{equation}
\label{eq:dft}
  X[k] = \sum_{n=0}^{N-1} x[n]\, e^{-j 2\pi k n / N}, \qquad k = 0,\ldots,N-1,
\end{equation}
em que \(X[k]\) é a componente espectral complexa do índice \(k\). A magnitude
\(|X[k]|\) indica a amplitude da componente na frequência \(k\,\Delta f\), e o
argumento \(\angle X[k]\) indica a sua fase. A Transformada Rápida de Fourier
(FFT) é um algoritmo eficiente para o cálculo da DFT, com complexidade
\(O(N \log N)\), viabilizando a análise espectral em tempo real ou quase real
\cite{cooley1965,oppenheim2010}.

A amplitude espectral de um pico corresponde, no domínio do tempo, à amplitude
de uma senoide de frequência igual à do pico. A fase, por sua vez, codifica a
posição angular da componente em relação a uma referência de tempo --- o
\emph{keyphasor} ou marcador de rotação --- e é essencial para distinguir
fenômenos como desbalanceamento estático e de acoplamento, que apresentam
relações de fase distintas entre mancais \cite{randall2011,harris2002}.

\section{O espaço \(R^3\) e as métricas do sinal}
\label{sec:r3}

A diretriz de persistência seletiva que fundamenta este trabalho propõe
representar cada componente espectral relevante como um ponto no espaço
tridimensional \(R^3\):
\begin{itemize}
  \item \textbf{frequência} --- identifica a origem cinemática do fenômeno
  (ordem \(kX\) em relação à rotação);
  \item \textbf{amplitude} --- quantifica a severidade do esforço mecânico;
  \item \textbf{fase} --- provê a correlação temporal para diagnósticos como
  desbalanceamento e análise orbital.
\end{itemize}

Além dos picos, o sistema calcula métricas de energia do sinal. O valor eficaz
(RMS) total de um sinal com \(N\) amostras é
\begin{equation}
\label{eq:rms}
  x_{RMS} = \sqrt{\frac{1}{N}\sum_{n=0}^{N-1} x[n]^2}.
\end{equation}
O RMS é a métrica padrão para a severidade global em aceleração e é a base de
critérios normativos como a ISO~10816 \cite{iso10816}. No contexto deste
trabalho, distinguem-se o RMS total, o RMS das componentes de pico (calculado a
partir das amplitudes dos picos \(R^3\)) e o RMS do ruído de fundo (obtido por
diferença de energia entre o total e os picos). O \emph{valor DC} é a média do
sinal e corresponde ao \emph{gap}/bias de sensores de proximidade (medição de
deslocamento), sendo persistido para permitir reconstrução de tendências.

\section{Catálogo de defeitos e assinaturas espectrais}

A análise de vibração relaciona padrões espectrais a defeitos mecânicos
específicos. O projeto adota um catálogo com treze tipos de defeito
(Tabela~\ref{tab:catalogo}), com assinaturas derivadas da literatura de
diagnóstico \cite{randall2011,scheffer2004} e formalizadas em um documento de
especificação de assinaturas utilizado como entrada para a geração sintética de
sinais.

\begin{table}[htbp]
\centering
\caption{Catálogo de tipos de defeito e assinatura espectral dominante.}
\label{tab:catalogo}
\begin{tabular}{@{}p{4.4cm}p{8.6cm}@{}}
\toprule
\textbf{Tipo} & \textbf{Assinatura dominante} \\
\midrule
sem\_defeito & 1X residual baixo, sem sequência forte de harmônicos \\
desbalanceamento & 1X dominante \\
desalinhamento\_angular & 2X + 1X, com 3X e fracionários 1,5X/2,5X \\
desalinhamento\_paralelo & 2X dominante + harmônicos pares 4X/6X/8X \\
rocamento & 1X + série de harmônicos; sub-harmônicos (0,5X, 1/3X, 2/3X) quando severo \\
mancal\_frouxo & família de harmônicos 1X--10X + fracionários \\
acoplamento\_defeituoso & 1X/2X + 3X/4X \\
oil\_whirl & subsíncrona 0,39X--0,48X (razão amostrada) + 1X \\
whirl\_atrito & subsíncrona não universal (faixa paramétrica, hipótese) \\
rolamento\_bpfo & BPFO + harmônicos $\pm$ sidebands 1X \\
rolamento\_bpfi & BPFI + harmônicos $\pm$ sidebands 1X \\
rolamento\_bsf & BSF + harmônicos \\
rolamento\_ftf & FTF + harmônicos \\
\bottomrule
\end{tabular}
\caption*{\footnotesize Fonte: SPECIFICATION.md do projeto; docs/assinaturas\_fft\_falhas\_maquinas\_rotativas\_IA.md.}
\end{table}

\subsection{Defeitos síncronos e harmônicos}

Forças centrífugas de \emph{desbalanceamento} crescem com o quadrado da
rotação e manifestam-se pela dominância de \(1X\) na direção radial; conforme o
tipo (estático, de acoplamento ou dinâmico), a fase entre mancais distingue o
quadro \cite{harris2002,randall2011}. O \emph{desalinhamento} angular produz
predominância de \(1X\) e \(2X\) na direção axial, com defasagem de \(180^\circ\)
através do acoplamento; o desalinhamento paralelo manifesta-se pela dominância
de \(2X\) na direção radial, com harmônicos pares superiores em casos severos.
\emph{Folgas mecânicas} (looseness) do tipo estrutural elevam \(1X\) vertical,
enquanto folgas móveis ou de mancal geram uma ``floresta'' de harmônicos
síncronos de \(1X\) até \(10X\). O \emph{roçamento} de rotor (rub) produz um
espectro rico em harmônicos e, distintamente, sub-harmônicos (\(1/2X\), \(1/3X\)),
com elevação do piso de ruído causada pelo truncamento da forma de onda
\cite{scheffer2004}.

\subsection{Instabilidades de filme de óleo}

Em mancais de deslizamento, a instabilidade de \emph{oil whirl} manifesta-se
por uma componente subsíncrona na faixa de aproximadamente \(0{,}40X\) a
\(0{,}49X\) da rotação, com órbita circular. Quando a frequência de whirl
coincide com uma frequência natural do sistema e ``trava'' nela, ocorre o
\emph{oil whip}, uma ressonância autoexcitada de caráter catastrófico
\cite{randall2011}. No catálogo do ARGUS, o \emph{oil whirl} é sintetizado com
razão amostrada uniformemente em \(0{,}39X\)--\(0{,}48X\) e o \emph{whirl por
atrito} como faixa paramétrica hipotética, sem razão universal.

\subsection{Falhas de rolamento}

Defeitos de mancais de rolamento geram frequências não síncronas características
que dependem da geometria do rolamento e da rotação. As frequências de falha
são, para um rolamento com \(n\) elementos rolantes, diâmetro do elemento
\(B_d\), diâmetro primitivo \(P_d\) e ângulo de contato \(\phi\)
\cite{randall2011,harris2002}:
\begin{align}
\label{eq:bpfo}
  BPFO &= \frac{n}{2}\, f_r \left(1 - \frac{B_d}{P_d}\cos\phi\right),\\[2pt]
\label{eq:bpfi}
  BPFI &= \frac{n}{2}\, f_r \left(1 + \frac{B_d}{P_d}\cos\phi\right),\\[2pt]
\label{eq:bsf}
  BSF  &= \frac{P_d}{2 B_d}\, f_r \left[1 - \left(\frac{B_d}{P_d}\cos\phi\right)^2\right],\\[2pt]
\label{eq:ftf}
  FTF  &= \frac{1}{2}\, f_r \left(1 - \frac{B_d}{P_d}\cos\phi\right),
\end{align}
em que BPFO corresponde a defeito na pista externa, BPFI na pista interna, BSF
nos elementos rolantes e FTF na gaiola. O surgimento de \emph{sidebands} em
torno dessas frequências indica modulação (tipicamente em \(1X\)) e severidade
avançada \cite{randall2011}. O ARGUS calcula essas ordens a partir de uma
geometria padrão (\(n=9\), \(B_d/P_d=0{,}2\), \(\phi=0\)), reproduzindo as
Equações~\eqref{eq:bpfo}--\eqref{eq:ftf} como múltiplos da frequência de
rotação; a geometria real do rolamento é uma pendência para uso em campo.

\subsection{Observação sobre amplitudes sintéticas}

As amplitudes relativas dos perfis adotados são \emph{parâmetros de síntese}
para geração de sinais e conjuntos de dados, e não valores universais de
severidade: a amplitude real depende de massa, rigidez, amortecimento,
posição do sensor, carga, ressonâncias e montagem. Essa distinção é essencial
para o uso academicamente honesto do simulador como ferramenta de treinamento,
e será retomada na discussão (Capítulo~\ref{cap:discussao}).

\section{Persistência seletiva e o critério \(R^3\)}

O segundo pilar do problema é o armazenamento de dados. A monitoração contínua
de turbomáquinas, com frequências de interesse de até 5~kHz, exige taxas de
amostragem entre 12,5~kHz e 50~kHz; multiplicadas por centenas de pontos de
medição, essas taxas geram um volume que inviabiliza o armazenamento
indiscriminado. A estratégia adotada --- e materializada no ARGUS --- é a
\emph{persistência seletiva}: migrar do domínio do tempo para o domínio da
frequência, conservando os vetores de picos \(R^3\) e as métricas de energia, e
descartar leituras redundantes cujos picos permaneçam dentro de uma zona de
tolerância em relação à última leitura efetivamente persistida.

A decisão de persistir ou descartar segue uma hierarquia lógica que evita o
\emph{drift} de dados (perda de pequenas variações acumuladas que seriam
ignoradas se a comparação fosse feita apenas contra a leitura anterior, que pode
ter sido descartada):

\begin{enumerate}
  \item \textbf{Variação de rotação} --- detecta mudança de regime operacional
  (rampas, carga). Pela \emph{regra de ouro}, a leitura é persistida se
  \begin{equation}
  \label{eq:regra-ouro}
    \left|RPM_{atual} - RPM_{armazenada}\right| > \Delta_{rot}
    \;\wedge\; (\text{número de picos} > 0).
  \end{equation}
  \item \textbf{Número de picos} --- identifica o surgimento de novas
  componentes espectrais.
  \item \textbf{Verificação \(R^3\)} --- aplica o \emph{Paralelepípedo de
  Descarte}: a leitura é descartada somente se, para cada pico atual, existir um
  pico de referência (o mais próximo em frequência, na última leitura
  persistida) tal que
  \begin{equation}
  \label{eq:descarte}
    |\Delta f| \le \delta_f \;\wedge\; |\Delta A| \le \delta_A
    \;\wedge\; |\Delta \phi| \le \delta_\phi,
  \end{equation}
  em que \(\delta_f\), \(\delta_A\) e \(\delta_\phi\) são as tolerâncias em
  frequência, amplitude e fase. Se ao menos um pico estiver fora dessa zona, a
  leitura é persistida e passa a ser a nova referência.
\end{enumerate}

A primeira leitura de um ponto é sempre persistida (condição de inicialização),
e a comparação é sempre feita contra a última leitura \emph{efetivamente
persistida}, nunca contra a última leitura avaliada/descartada. As leituras
descartadas são registradas em um armazenamento separado (\emph{trash}),
permitindo auditar a taxa de descarte e validar o algoritmo. Esse critério
viabiliza a redução do volume de dados sem perda da capacidade diagnóstica,
direcionando o esforço analítico para as transições de estado relevantes do
ativo \cite{scheffer2004}.

\section{Considerações normativas}

O trabalho observa, na sua apresentação, as normas ABNT de documentação
(NBR~14724 e NBR~6023) \cite{abnt14724,abnt6023} e, na fundamentação técnica,
as normas ISO de avaliação de vibração e qualificação de pessoal
\cite{iso10816,iso18436}. As edições vigentes dessas normas na data da defesa
devem ser verificadas, conforme registrado nas pendências do projeto.
